\documentclass[10pt,a4paper]{amsart}
\usepackage[margin=19mm]{geometry}
\usepackage{amsmath,amssymb,mathtools}
\usepackage[scr=rsfs,cal=euler]{mathalfa}
\usepackage{xfrac}
\usepackage[unicode,hidelinks,bookmarks=false]{hyperref}
\newcommand{\ie}{i.e.\ }
\newcommand{\cf}{cf.\ }
\newcommand{\resp}{respectively\ }
\newcommand{\Subsection}{\subsection}
\providecommand{\nobreakdash}{\nobreak\hbox{-}}
\setlength{\emergencystretch}{2em}
\multlinegap0pt
\usepackage{bm}
\usepackage{bbm}
\usepackage{dsfont}
\usepackage{xspace}
\usepackage{cancel}
\usepackage{mathtools}
\usepackage{amsthm}
\makeatletter
\@ifundefined{theorem}{\newtheorem{theorem}{Theorem}}{}
\@ifundefined{lemma}{\newtheorem{lemma}[theorem]{Lemma}}{}
\makeatother
\newcommand{\pHz}[2]{\mathcal{H}^{#1}(#2)}
\title[On sums involving powers of harmonic numbers]{On sums involving powers of harmonic numbers}
\author{Lo Ho Tin}
\date{Source: arXiv:2507.03058v1 (CC BY 4.0)}
\begin{document}
\maketitle

\noindent\emph{Source and licence.} On sums involving powers of harmonic numbers. Version arXiv:2507.03058v1 (CC BY 4.0), distributed under the Creative Commons Attribution 4.0 International licence, as recorded in the arXiv licence metadata for this exact version and shown on the abstract page https://arxiv.org/abs/2507.03058v1. Full rights notice: licenses/arxiv-2507.03058-CC-BY-4.0.txt.\par\smallskip

\noindent\textit{Editorial scope.} Counted unit: the Theorem whose source label is \texttt{Laurent expansion of H} in the article (source0.tex lines 638--644, source label \texttt{Laurent expansion of H}), delivered together with the article's own complete original proof of it (source0.tex lines 645--669, one \texttt{proof} environment of 25 lines). No mathematics is written, repaired or re-derived: statement and proof are the article's own text line for line. Required context retained uncounted: Asymp. Harmonic (lines 119--119); formal power series (lines 623--629). Every definition, display and label of those lines is retained unchanged. Omitted from this package and not redistributed: 1--118, 120--622, 630--637, 670--694. Nothing inside a retained excerpt is omitted or altered: every retained body is delivered exactly as the article's lines are written, so no omission receipt of any kind accompanies this package. Citations: the retained excerpts carry no citation command at all, so no bibliography entry of the article is transcribed and no reference is redistributed. Reference adaptation: none was needed and none was made; every reference inside the delivered unit resolves inside it, so no unresolved reference or citation is produced. Printed slips kept literal: no printed word, symbol, hypothesis or number of the counted statement or of its proof was altered. Layout and class: the article's own class is \texttt{article} with packages including multirow, amsmath, amssymb, amsfonts, amsthm, mathrsfs, appendix, xcolor, textcomp, manyfoot, booktabs, algorithm, algorithmicx, algpseudocode; the wrapper is the lane's pinned \texttt{amsart} editorial wrapper at 10pt on A4, which restates the theorem-like environments the retained text needs and adds the article's own macro definitions that the retained text needs (\texttt{\textbackslash pHz}), with extra wrapper packages (bm, bbm, dsfont, xspace, cancel, mathtools). The delivered numbering is the unit's own. Assets: no retained span contains a figure environment, a graphics-inclusion command, a PDF-page inclusion, a table or a TikZ picture; no image file is copied into this package, the package declares no asset dependency and no image file is redistributed with it. Licence: version arXiv:2507.03058v1 of this article is distributed under the Creative Commons Attribution 4.0 International licence, as recorded in the arXiv licence metadata for this exact version and shown on the abstract page https://arxiv.org/abs/2507.03058v1. A full rights notice is delivered at licenses/arxiv-2507.03058-CC-BY-4.0.txt. Characters: every retained excerpt is pure ASCII, so the delivered engine, XeLaTeX with its default Latin Modern text font, reports no missing character.
\begin{equation}\label{Asymp. Harmonic}H_n=\log(n)+\gamma +\frac{1}{2n}-\sum_{a=1}^k \frac{B_{2a}}{2an^{2a}} +\int_n^{\infty} \frac{\tilde B_{2k}(x)}{x^{2k+1}}\, dx\end{equation}

\begin{lemma}\label{formal power series}
    Consider the formal power series \[f(z)=\sum_{n=1}^{\infty} (-1)^n \zeta(1-n)z^n\]
    Then \begin{equation}\label{generating function of a}
        [f(z)]^m =\left(\frac{z}{2}+\zeta(-1)z^2+\zeta(-3)z^4+\zeta(-5)z^6+\cdots\right)^m=\sum_{n\geq m} a_n(m)z^n
    \end{equation}
and \[(\gamma+ f(x))^m=\sum_{n=0}^{\infty} H^m(1-n) x^n\]
\end{lemma}

\begin{theorem}\label{Laurent expansion of H}
    Let $\mathrm{H}^m(n,k)$ denote the $(s-n)^{-k}$ coefficient of the Laurent expansion of $\pHz{m}{s}$ expanded at $s=n$. Then
    \begin{equation}\label{coefficient of expansion}\mathrm{H}^m(1-n,k)=\sum_{\ell=0}^m \binom{m}{\ell}\gamma^{m-\ell} a_n(\ell,k)\end{equation}
    Where $a_n(\ell, k)$ is a triple sequence of rational numbers and takes the form $\ell(\ell-1)\cdots(\ell -k+2)a_n(\ell-k+1)$. Moreover, we have the formal power series \begin{equation}
        \sum_{n=0}^{\infty}\mathrm{H}^m(1-n,k+1)x^n=m(m-1)\cdots(m-k+1)(\gamma+f(x))^{m-k}
    \end{equation}
\end{theorem}

\begin{proof}
    We will use the symbol $\sum_{n=0}$ to denote the sum has an upper bound that is yet determined. Using (\ref{Asymp. Harmonic}) we get\begin{equation}\label{relating m+1 to m}
        \mathcal{H}^{m+1}(s)=-{\mathcal{H}^m}'(s)+\gamma \mathcal{H}^m (s)+\sum_{a=1}(-1)^a \zeta(1-a)\mathcal{H}^m(a+s)+\cdots
    \end{equation} Pulling out the terms gives
    \[\mathrm{H}^{m+1}(n,k)=(k-1)\mathrm{H}^m(n,k-1)+\gamma\mathrm{H}^m(n,k)+\sum_{a=1}(-1)^a \zeta(1-a)\mathrm{H}^m(n,k)\]
If we set $n\mapsto 1-n$ and write H in the form (\ref{coefficient of expansion}), which we know is true for $k=1,2,3$ and is easy to be seen true in general, we get
\[\begin{split}
   & \mathrm{H}^{m+1}(1-n,k)\gamma^{-m}- \mathrm{H}^m(1-n,k)\gamma^{1-m} \\ & =(k-1)\sum_{l=0}^{\infty}\binom{m}{l}\frac{1}{\gamma^l}a_n(l,k-1) +\sum_{a=1} (-1)^a\zeta(1-a)\sum_{l=0}^{\infty}\binom{m}{l}\frac{1}{\gamma^l}a_{n-a}(l,k)
\end{split}\]
summing both sides in $m$ gives a telescoping sum, which is simplified as
\[\mathrm{H}^m(n,k)\gamma^{-m}=\sum_{l=1}^{\infty}\binom{m}{l}\frac{1}{\gamma^{l}}\left\{(k-1)a_n(l-1,k-1)+\sum_{a=1}(-1)^a \zeta(1-a)a_{n-a}(l-1,k)\right\}\]
Equating coefficients gives
\[a_n(l,k)=(k-1)a_n(l-1,k-1)+\sum_{a=1}(-1)^a \zeta(1-a)a_{n-a}(l-1,k)\]
Consider the generating function of both sides. For the sake of induction, let's assume that $a_n(l,k-1)$ could be written in the form $l(l-1)\cdots(l-k+3)a_n(l-k+2)$, which we proved in the prior section for $k=2,3,4$. Set the formal power series $f_{k,l}(x)=a_1(l,k)x +a_2(l,k)x^2+a_3(l,k)x^3+\cdots$. Then we can write the above as 

$$f_{k,l}(x)=(k-1)f_{k-1,l-1}(x)+f(x)f_{k,l-1}(x)$$
I may omit the $x$ and write $f(x)$ as $f$ from now on. Using lemma \ref{formal power series}, we have $f_{k-1,l-1}=(l-1)\cdots(l-k+2)f^{l-k+1}$. Therefore

\[\begin{split}
    f_{k,l}-f\cdot f_{k,l-1}&=(k-1)(l-1)\cdots(l-k+2)f^{l-k+1}\\ 
\frac{f_{k,l}}{f^l}-\frac{f_{k,l-1}}{f^{l-1}}& =(k-1)f^{1-k}\cdot (l-1)\cdots (l-k+2)
\end{split}\]
Summing both sides gives
\[\frac{f_{k,m}}{f^m}-\frac{f_{k,0}}{f^0}=\frac{m(m-1)\cdots(m-k+2)}{f^{k-1}}\]
Note that $f_{k,0}=0$ for $k\geq 2$, Finally we multiply both sides by $f^m$ and equate coefficients to show induction.\end{proof}

\end{document}
